Giao dịch tiền mã hóa đòi hỏi kết hợp dữ liệu giá, dữ liệu giao dịch trên chuỗi khối (on-chain) và tin tức, vốn
thường đưa ra tín hiệu trái chiều. Các agent giao dịch dựa trên mô hình ngôn ngữ lớn (LLM) hiện nay, tiêu biểu là
CryptoTrade (EMNLP 2024), gộp các nguồn trong một lời nhắc và để mô hình tự xử lý mâu thuẫn một cách ngầm định.

Đề tài xây dựng một hệ thống đa tác tử trong đó mỗi nguồn dữ liệu do một agent chuyên trách phân tích; mâu thuẫn giữa
các agent được đo tường minh bằng chỉ số kết hợp phân kỳ Kullback--Leibler và phương sai quyết định; khi mâu thuẫn vượt
ngưỡng, hệ thống phân tích nguyên nhân gốc và tổ chức tranh biện nhiều vòng, trong đó mọi lập luận phải trích dẫn
bằng chứng gốc và được kiểm tra; cuối cùng một Mediator tổng hợp các nhận định thành tín hiệu BUY/SELL/NEUTRAL.

Đề tài cài đặt lại CryptoTrade, kiểm chứng khớp với mã nguồn gốc, rồi chạy cả hai hệ thống trên cùng dữ liệu, cùng môi
trường giao dịch, cùng mô hình ngôn ngữ chạy cục bộ, trên ba giai đoạn thị trường năm 2023 của bài báo và một tập phát
triển gồm 145 ngày khác. Trên BTC, hệ thống đề xuất đạt +0,38\% trên tập kiểm thử và +8,60\% trên tập phát triển, thấp
hơn CryptoTrade đầy đủ (+5,71\% và +11,95\%). Phân tích cho thấy: (i) tranh biện giảm mạnh mức mâu thuẫn (từ 0,61 xuống
0,22) nhưng chỉ thay đổi 11 trên 195 quyết định, và ở các ngày agent mâu thuẫn nó phân xử đúng 48,3\%, ngang việc không
tranh biện; (ii) từng agent đoán hướng gần như ngẫu nhiên ở tầm 1--7 ngày và bị lệch phe: Sentiment Agent nói BUY khoảng
80\% số lần, còn Market Agent bị các nhãn chỉ báo đảo chiều trong dữ liệu kéo về SELL giữa các đợt tăng; (iii) tám hướng
cải tiến, chọn trên tập phát triển và đánh giá một lần trên tập kiểm thử, không vượt được CryptoTrade một cách ổn định,
còn bỏ các lệnh SELL đưa lợi suất lên ngang CryptoTrade; (iv) CryptoTrade cũng không đoán hướng tốt hơn ngẫu nhiên mà
mua ở 91\% số ngày có hướng. Một kiểm chứng đăng ký trước trên ETH xác nhận các phát hiện này, bác bỏ một cải tiến được
tìm ra sau khi xem dữ liệu (trọng tài theo độ tin cậy: tỷ lệ đúng giảm từ 52,3\% xuống 46,9\%), và cho thấy thứ hạng
đảo chiều: trên ETH hệ thống đề xuất cao hơn CryptoTrade 7,85 điểm trên tập kiểm thử. Không chênh lệch lợi suất nào giữa
hai hệ thống có ý nghĩa thống kê. Kết quả cho thấy hạn chế chính nằm ở chất lượng nhận định của từng agent chứ không ở
cách tổng hợp, và lợi suất của các agent giao dịch LLM phản ánh mức tiếp xúc thị trường nhiều hơn khả năng dự đoán.

\paragraph{Từ khóa:} hệ đa tác tử, mô hình ngôn ngữ lớn, tranh biện đa tác tử, phát hiện mâu thuẫn, giao dịch tiền mã hóa.
